% Micfong Sans Math in a LaTeX document, with IBM Plex for the text.
% From this folder: xelatex latex.tex
% The math setup works the same under LuaLaTeX. Only the CJK paragraphs depend on
% xeCJK, which is XeLaTeX-only, and they need IBM Plex Sans SC, TC, JP, KR and
% Devanagari installed. IBM Plex Sans and Mono come with TeX Live.
\documentclass[11pt]{article}
\usepackage[a4paper, margin=2.4cm]{geometry}
\usepackage{amsmath}
\usepackage{unicode-math}
\usepackage{upgreek}

\setmainfont{IBMPlexSans}[
  Extension = .otf, UprightFont = *-Regular, ItalicFont = *-Italic,
  BoldFont = *-SemiBold, BoldItalicFont = *-SemiBoldItalic]
\setmonofont{IBMPlexMono}[Extension = .otf, UprightFont = *-Regular, Scale = 0.88]

% The math font. Path points at the fonts in this repository; drop it once the
% fonts are installed. BoldFont gives \boldmath its own weight.
\setmathfont{MicfongSansMath.otf}[
  Path = ../dist/, BoldFont = MicfongSansMath-Bold.otf]

\usepackage{xeCJK}
\xeCJKsetup{CJKspace = true}
\setCJKmainfont{IBM Plex Sans SC}
\setCJKmonofont{IBM Plex Sans SC}
\newCJKfontfamily\tcfont{IBM Plex Sans TC}
\newCJKfontfamily\jpfont{IBM Plex Sans JP}
\newCJKfontfamily\krfont{IBM Plex Sans KR}
\newfontfamily\devanagarifont{IBM Plex Sans Devanagari}[Script = Devanagari]

\usepackage{xcolor}
\usepackage{titlesec}
\titleformat*{\section}{\large\bfseries}
\setlength{\parindent}{0pt}
\setlength{\parskip}{0.5em}
% \lang[font switch]{language name}{paragraph}
\newcommand{\lang}[3][]{\par\noindent
  \makebox[2.4cm][l]{#1\footnotesize\color{gray}#2}%
  \parbox[t]{\dimexpr\linewidth-2.4cm}{#1\raggedright #3}\par\medskip}

\begin{document}

\begin{center}
  {\LARGE\bfseries Micfong Sans Math}\par\smallskip
\end{center}

The equations in this document are set in Micfong Sans Math, a new open-source font for mathematics. The font is designed to be compatible with IBM Plex Sans, which is used for the body text.

To use Micfong Sans Math in your own LaTeX documents, install the font and \texttt{unicode-math}, then add the following line to your document:

\begin{verbatim}
\setmathfont{MicfongSansMath.otf}[BoldFont = MicfongSansMath-Bold.otf]
\end{verbatim}

\section{Fourier series}

Let $f\colon \mathbb{R} \to \mathbb{C}$ be $2\uppi$-periodic and integrable on $[-\uppi, \uppi]$.
Its Fourier coefficients are
\begin{equation}
  \hat{f}(n) = \frac{1}{2\uppi} \int_{-\uppi}^{\uppi} f(x) \mathrm{e}^{-\mathrm{i}nx} \,\mathrm{d}x, \quad n \in \mathbb{Z},
\end{equation}
and Parseval's identity states that
\begin{equation}\label{eq:parseval}
  \sum_{n \in \mathbb{Z}} \lvert\hat{f}(n)\rvert^2 = \frac{1}{2\uppi} \int_{-\uppi}^{\uppi} \lvert f(x)\rvert^2 \,\mathrm{d}x.
\end{equation}
For $f(x) = x$ the coefficients are $\hat{f}(0) = 0$ and $\hat{f}(n) = (-1)^n \mathrm{i}/n$, so
\eqref{eq:parseval} gives $2\sum_{n \ge 1} 1/n^2 = \uppi^2/3$, which is Euler's $\zeta(2) = \uppi^2/6$.

\section{Linear algebra}

A real symmetric matrix $\mathbf{A} \in \mathbb{R}^{n \times n}$ has an orthonormal basis of eigenvectors,
so $\mathbf{A} = \mathbf{Q} \boldsymbol{\Lambda} \mathbf{Q}^\top$ with $\mathbf{Q}^\top \mathbf{Q} = \mathbf{I}$ and $\boldsymbol{\Lambda} = \operatorname{diag}(\lambda_1, \dots, \lambda_n)$.
For example,
\begin{equation}
  \begin{pmatrix} 2 & 1 \\ 1 & 2 \end{pmatrix}
  = \frac{1}{2} \begin{pmatrix} 1 & -1 \\ 1 & 1 \end{pmatrix}
  \begin{pmatrix} 1 & 0 \\ 0 & 3 \end{pmatrix}
  \begin{pmatrix} 1 & 1 \\ -1 & 1 \end{pmatrix}, \quad
  \det(\mathbf{A} - \lambda \mathbf{I}) = (\lambda - 1)(\lambda - 3).
\end{equation}
The Rayleigh quotient $R(\boldsymbol{x}) = (\boldsymbol{x}^\top \mathbf{A} \boldsymbol{x}) / (\boldsymbol{x}^\top \boldsymbol{x})$
takes values in $[\lambda_{\min}, \lambda_{\max}]$.

\section{Probability}

A random variable $X \sim \mathcal{N}(\mu, \sigma^2)$ has density
\begin{equation}
  \varphi(x) = \frac{1}{\sigma\sqrt{2\uppi}} \exp\biggl(-\frac{(x - \mu)^2}{2\sigma^2}\biggr),
\end{equation}
with $\mathbb{E}[X] = \mu$ and $\operatorname{Var}(X) = \sigma^2$. Whatever the distribution,
as long as the variance is finite, Chebyshev's inequality bounds the tails:
\begin{equation}
  \mathbb{P}(\lvert X - \mu\rvert \ge k\sigma) \le \frac{1}{k^2}, \quad k > 0.
\end{equation}

\section{Physics}

Maxwell's equations in a vacuum:
\begin{equation}
\begin{aligned}
  \nabla \cdot \boldsymbol{E} &= \frac{\rho}{\varepsilon_0}, &
  \nabla \times \boldsymbol{E} &= -\frac{\partial \boldsymbol{B}}{\partial t}, \\
  \nabla \cdot \boldsymbol{B} &= 0, &
  \nabla \times \boldsymbol{B} &= \mu_0 \boldsymbol{J} + \mu_0 \varepsilon_0 \frac{\partial \boldsymbol{E}}{\partial t}.
\end{aligned}
\end{equation}
The time-dependent Schrödinger equation is
$\mathrm{i}\hbar\, \partial_t \Psi(\boldsymbol{r}, t) = \hat{H} \Psi(\boldsymbol{r}, t)$ with
$\hat{H} = -\frac{\hbar^2}{2m} \nabla^2 + V(\boldsymbol{r})$.

\section{Multilingual support}

Micfong Sans Math takes its letters from IBM Plex Sans, so it sits naturally beside the other
IBM Plex families. Each paragraph below is set in the matching one.

\lang{简体中文}{设 $f$ 在区间 $[a, b]$ 上连续，且 $F' = f$，则
  $\int_a^b f(x) \,\mathrm{d}x = F(b) - F(a)$。这就是微积分基本定理。}
\lang[\tcfont]{繁體中文}{若 $\lvert q\rvert < 1$，則幾何級數 $\sum_{n=0}^{\infty} q^n$ 收斂，且其和為 $1/(1 - q)$。}
\lang[\jpfont]{日本語}{オイラーの公式 $\mathrm{e}^{\mathrm{i}\theta} = \cos\theta + \mathrm{i} \sin\theta$ において $\theta = \uppi$ とおくと、オイラーの等式 \mbox{$\mathrm{e}^{\mathrm{i}\uppi} + 1 = 0$} が得られる。}
\lang[\krfont\xeCJKsetup{CJKecglue = {}}]{한국어}{직각삼각형에서 빗변의 길이를 $c$, 나머지 두 변의 길이를 $a$, $b$라고 하면 $a^2 + b^2 = c^2$이 성립한다.}
\lang{Русский}{Для любого $\varepsilon > 0$ существует $\delta > 0$ такое, что
  $\lvert f(x) - f(a)\rvert < \varepsilon$ для всех $x$, удовлетворяющих $\lvert x - a\rvert < \delta$.
  Тогда говорят, что функция $f$ непрерывна в точке $a$.}
\lang{Ελληνικά}{Το εμβαδόν κυκλικού δίσκου ακτίνας $\rho$ είναι $E = \uppi \rho^2$, ενώ το μήκος του κύκλου είναι $L = 2\uppi\rho$.}
\lang[\devanagarifont]{हिन्दी}{द्विघात समीकरण $ax^2 + bx + c = 0$, जहाँ $a \ne 0$, के मूल
  $x = \frac{-b \pm \sqrt{b^2 - 4ac}}{2a}$ होते हैं।}

\end{document}
